\section{Discussion}\label{sec:discussion}

\paragraph{Advanced control buys comfort first, and savings second.}
Against a realistic split-AC thermostat, none of the modulating controllers saves money
in this dwelling. The thermostat's low cost is not efficiency: it comes from short,
full-capacity bursts that overshoot below the comfort band (\KhColdTh\,K\,h of
cooling-induced overcooling) and from about \StartsRatioTh{} times more compressor
starts, and its 15-min peak demand is 21--26\,\% higher than that of PID, E-MPC and
Hybrid. Studies that report large savings against weak or non-modulating baselines
should therefore be read with care. The original version of this study is an example:
its rule-based baseline cooled regardless of indoor temperature. The fair comparison is
between controllers that deliver similar comfort. On that basis the Hybrid controller
is \RelHybPidCost\,\% cheaper than SIMC-tuned PID and has \RelHybPidKh\,\% less
discomfort. Both differences are statistically significant, and both hold under
heatwave, thermal-mass and tariff variations.

\paragraph{Why the Hybrid controller helps.}
The economic MPC finds the cheap operating point: it keeps occupied zones just inside
the upper comfort bound and pre-cools the slab at night when that pays. It does so with
hourly decisions from a linear model that is about 10--15\,\% wrong and from imperfect
forecasts. The ablations show that the Hybrid's two lower layers do different jobs.
\begin{itemize}
  \item The 15-min PI layer supplies essentially all of the cost saving over E-MPC.
  \item Most of the comfort gain comes from the E-MPC with the safety override alone,
        which forces near-full capacity in an occupied zone that exceeds the band.
\end{itemize}
The perfect-information ablation is consistent with this: a better model and perfect
weather do not close the gap, so the gap comes from features outside the MPC's linear
model and from its choice to run close to the bound. The heatwave back-off test points
the same way. Comfort under sustained heat is governed mainly by how much margin the
controller keeps, not by how well it predicts.

\paragraph{Little ToD arbitrage.}
The ToD tariff barely moves the load in this dwelling, even under an economic MPC with a
24-h horizon:
\begin{itemize}
  \item every kWh of peak-window cooling serves an occupied zone, so it cannot be
        dropped without discomfort (occupied zone-hours fall in the peak windows
        slightly less often than hours in general: 27\,\% against 33\,\%);
  \item the heat stored in the roof slab is released in the afternoon and evening:
        mean slab-to-air heat flow is about 2.5\,kW during 18:00--23:00, against
        1.8\,kW at 10:00--16:00;
  \item the peak-to-off-peak price ratio (1.7) is modest. At 3.6:1 the predictive
        controllers move about 5 percentage points of cooling energy out of the peak
        windows.
\end{itemize}
The night pre-cooling we observe mainly reduces late-morning demand, which falls in a
normal-tariff window. Larger demand shifting would need steeper tariffs than the one tested here,
thermal storage or a roof with more accessible mass. The original study reported
evidence of tariff-driven shifting; with real hourly profiles and costs, that evidence
does not survive.

\paragraph{Reinforcement learning needs engineering to be competitive.}
An unconstrained DQN learned to cool through the winter, because "do nothing" is
rarely rewarded once weather-driven cold is excluded from the penalty. That agent was
uncompetitive. Two changes made the DQN competitive with E-MPC on comfort and with
PID-to-MPC on cost:
\begin{itemize}
  \item a safety shield, with the agent's own proposals stored for learning;
  \item checkpoint selection on a separate validation set.
\end{itemize}
The DQN also needs no model at run time. Its seed-to-seed spread is an order of
magnitude larger than that of the model-based controllers, so a single training run is
not a reliable estimate of performance. This supports the protocol used here: several
seeds, a validation set separate from the test year, and robustness checks on weather
not seen in training.

\paragraph{IAQ is decided by design, not by the thermostat logic.}
Indoor CO$_2$ and PM$_{2.5}$ are essentially identical across thermal controllers. They
are governed by the ventilation rate and the filtration architecture: the recirculating
HEPA filter alone reduces mean PM$_{2.5}$ more than threefold. Humid hours are where
design and control interact. The demand-controlled dehumidifying DOAS cuts occupied
hours above 70\,\% RH by about a factor of seven compared with no mechanical
ventilation, and the steadier modulating controllers remove more moisture than the
thermostat. The DOAS (fans and coil) uses about a quarter of annual electricity in this
climate. That load deserves its own optimisation, for example demand-controlled
ventilation or a combined IAQ-thermal MPC.

\paragraph{Practical implications.}
For Indian homes with inverter VRF or multi-split systems:
\begin{enumerate}
  \item Modulating control in place of thermostat on/off is the largest gain. It
        delivers comfort and fewer compressor starts, and with PID or the model-based
        controllers about 18\,\% lower peak demand.
  \item A hybrid MPC plus feedback architecture adds a further 3--4\,\% cost reduction
        over well-tuned PID at better comfort.
  \item Ventilation and filtration design decide IAQ and should be treated as a design
        variable of equal standing.
\end{enumerate}
The computational cost of the MPC (\MsMpc\,ms per hourly solve in interpreted
MATLAB/Octave) is compatible with an embedded gateway.
